\subsection*{The discrete source groups remain disjoint under every conjugation}

\statement{B.72}{Proved}{Every discrete mixed conjugate intersection is trivial}
For $\Gamma=\Gamma_{0,[12]}$, and the actual subgroups
$H_T=C_P(T)$ and $H_5=C_P(D)$ of B.28, one has
\[
                  H_T\cap gH_5g^{-1}=1
                       \qquad\text{for every }g\in\Gamma.
                                                               \tag{B.72.1}
\]
This is a statement in the discrete mapping-class group.
Neither the corresponding intersection of completed subgroups
nor either factor's individual malnormality is asserted.

\begin{proof}
\emph{The four-point comparison and the cyclic Nielsen--Thurston types.}
In B.28's cyclic coordinates the covering is
\[
 z\longmapsto z^5:\quad
 \mathbb P^1\setminus\{0,\infty,\mu_5,u\mu_5\}
 \longrightarrow
 \mathbb P^1\setminus\{0,\infty,1,u^5\}.
                                                               \tag{B.72.2}
\]
It is an unramified covering of punctured surfaces of degree five:
both ramification poles have been removed.
The base moduli parameter is $v=u^5$, and
\[
 \mathbb P^1_u\setminus\{0,\infty,\mu_5\}
      \longrightarrow\mathbb P^1_v\setminus\{0,1,\infty\}
                                                               \tag{B.72.3}
\]
is an unramified degree-five cover as well.
The faithful quotient comparison in B.28 identifies $H_5$
with its index-five subgroup of $P\Gamma_{0,4}\simeq F_2$.
The deck group of (B.72.2) has no nonidentity pure element,
so this comparison has no pure finite kernel.

A nonidentity pure four-point mapping class is either
pseudo-Anosov or a nonzero power of a peripheral twist.
For example this is the hyperbolic/parabolic classification
in its torsion-free finite-index modular-group model;
the three peripheral classes are the three punctures of
the target in (B.72.3).
Every lift of a pseudo-Anosov map in (B.72.2) is again
pseudo-Anosov: pull back its two transverse measured
foliations and their reciprocal dilation factors.
All covering ramification is at marked punctures, so
there is no filling step that might remove a one-prong
singularity or alter that conclusion.

The inverse images of the three peripheral classes in
(B.72.3) are exactly the seven source cusps.
Their lifts have the following explicit reduction types:
\[
\begin{array}{c|c|c}
\text{source cusp}&\text{colliding marked clusters}
                       &\text{number of twist curves}\\\hline
0,\infty&\text{one pole and one orbit of five marks}&1\\
u\in\mu_5&\text{five disjoint pairs}&5.
\end{array}                                                     \tag{B.72.4}
\]
At zero the shrinking orbit and its fixed pole form a
six-point cluster; at infinity use $u^{-1}$ and rescale.
At a fifth root of unity, five different pairs have
simple transverse collisions in disjoint small disks.
A positive parameter meridian therefore gives a full
twist about every listed cluster boundary, with a
nonzero coefficient at each boundary.
The boundaries are essential and pairwise nonisotopic.
A nonzero power has the same canonical reduction
system, consisting of exactly these one or five curves.
This describes every reducible nonidentity element
of $H_5$, up to conjugacy in $H_5$.

Now suppose $1\ne h\in H_T\cap gH_5g^{-1}$.
Then $h$ centralizes both $T$ and $D'=gDg^{-1}$.

\emph{The reducible case.}
Conjugacy preserves (B.72.4)'s numbers of reduction curves.
Since $T$ centralizes $h$, it acts on that canonical
reduction system, a set of size one or five.
Every such $A_4$-action has a fixed point.
Indeed an orbit size divides twelve, and the only
possible orbit sizes at most five, apart from one,
are three and four: $A_4$ has no index-two subgroup.
Neither a sum of threes and fours nor a single such
number equals five.

Let $c$ be a reduction curve whose isotopy class is
fixed by $T$.  Its two sides define an unordered
partition of the twelve punctures.
The induced action on the two parts is a homomorphism
$T\to C_2$, necessarily trivial because
$A_4^{\mathrm{ab}}\simeq C_3$.
Both parts would therefore be $T$-invariant.
An essential curve has nonempty proper parts,
contradicting the actual transitive $T$-action on
the twelve punctures.  Thus $h$ cannot be reducible.

\emph{The pseudo-Anosov case and its finite symmetry kernel.}
We recall the precise finite-kernel argument needed here,
including its marked-sphere scope.
For a pseudo-Anosov mapping class, a commuting mapping
class has a representative preserving each of the two
invariant foliations, multiplying their transverse
measures by positive reciprocal constants.
The usual uniqueness of the pseudo-Anosov representative
is applied relative to the marked punctures.
Multiplication of these constants gives a homomorphism
from the centralizer to $\mathbb R_{>0}$.
Its kernel is finite.

Here is the latter assertion directly in the punctured
case.  Include the separatrices at one-prong marked
punctures as well as those at ordinary singularities.
There are finitely many and at least one on this
negative-Euler-characteristic marked sphere.
A symmetry preserving both transverse measures and
fixing one separatrix fixes its endpoint and its
transverse-measure length coordinate, hence fixes
that dense leaf pointwise and is the identity.
The symmetry kernel therefore acts freely on a finite
set and is finite.
The foliation-preserving representatives inject into
the mapping-class group: one representing the identity
relative to the punctures has dilation one, belongs
to that finite symmetry group, and is periodic while
fixing all twelve punctures.  An orientation-preserving
periodic sphere homeomorphism fixing at least three
points is the identity.
These are the steps of the dilation and finite-symmetry
proof of \citet[Lemmas 1--4]{DE_McCarthy1994}, with marked
separatrices and relative isotopies retained explicitly;
we do not substitute the punctured case silently for
the compact negative-Euler hypothesis in its introduction.

Every finite subgroup of the centralizer lies in this
finite kernel, since $\mathbb R_{>0}$ has no nontrivial
finite subgroup.  Consequently
\[
                         E=\langle T,D'\rangle
                                                               \tag{B.72.5}
\]
is itself a finite subgroup of $\Gamma$.
Nielsen realization and the finite rotation-group
classification show that $E\simeq A_5$.
Indeed cyclic and dihedral rotation groups contain
no $A_4$, and neither $A_4$ nor $S_4$ has an element
of order five.  These exhaust the alternatives to $A_5$.

The twelve punctures must form the unique exceptional
$E$-orbit of size twelve: its other nontrivial orbit
sizes are twenty, thirty and sixty.
Thus this is an actual rigid icosahedral subgroup,
conjugate to the one in B.7/B.52.
Its discrete centralizer is trivial by the quotient
triangle argument there.
But $h$ centralizes every generator of (B.72.5), hence
centralizes $E$, a final contradiction.
\end{proof}

The proof supplies a geometric distinction between the
two source subgroups throughout the discrete ambient group.
It does not make a profinite element pseudo-Anosov or
peripheral, and it proves no density assertion for their
completed intersection.  In particular the conclusion
does not remove B.46's finite-index refinement condition
and leaves $\mathcal K$ open.
